\subsection{A nonfree differential image removes the first horizontal obstruction}

\begin{proposition}[D-027: \status{Proved}; the first row of a standard-tail image]
\label{sa08:D027}
Let \(k\) be any commutative ring, \(S=\Gamma_k(T)\), and let
\(F\) be the standard simplicial PD de Rham theory. Fix \(N\ge2\).
Define \(D_N(k)\subset F\) to be the differential PD subalgebra
generated by \(S\) and all the degree-zero elements of weights
at least \(N\). Equivalently, its degree-zero algebra is
\[
 D_N^{w,0}=
 \begin{cases}S^w,&w<N,\\ F^{w,0},&w\ge N,\end{cases}
\]
and it is closed under the differential, products and the
defined divided powers in \(F\).
Then, for its simplicial form-degree-one row,
\[
 \pi_0(D_N^{w,1})=0,\qquad
 \pi_1(D_N^{w,1})=
 \begin{cases}
 k,&N\le w<2N,\\
 0,&w<N\text{ or }w\ge2N.
 \end{cases}
 \tag{D-27a}
\]
Moreover, in every simplicial degree, the vertical cohomology
of \(D_N\) is \(S\) in degree zero and zero in degree one.
These assertions do not assert vanishing of all higher
horizontal homotopy groups or all higher vertical cohomology.
\end{proposition}

\begin{proof}
In form degree one, the subalgebra consists exactly of finite
sums \(a\,df\), where \(f\) has weight at least \(N\) and \(a\)
belongs to \(D_N^0\). Higher divided powers and products of
positive-degree forms cannot add further form-degree-one
generators. A coefficient \(a\in S\) gives \(a\,df=d(af)\).
For \(w\ge N\), let
\[
 Z_w=d(F^{w,0})\subset D_N^{w,1}.
\]
The simplexwise Poincar\'e exactness gives a short exact
sequence of simplicial modules
\[
 0\longrightarrow S^w\longrightarrow F^{w,0}
             \xrightarrow{d} Z_w\longrightarrow0.
 \tag{D-27b}
\]
The positive-weight simplicial module \(F^{w,0}=\Gamma^w(\Phi)\)
is contractible, while \(S^w\cong k\) is constant. Thus
\(\pi_1(Z_w)\cong k\) and \(\pi_j(Z_w)=0\) for \(j\ne1\).

There are no one-forms in simplicial degree zero. In degree
one use coordinates \(T,X\). Every basis form
\(T^{[c]}X^{[d]}dX\), of weight \(c+d+1=w\), is the
differential of \(T^{[c]}X^{[d+1]}\).
Consequently
\[
              (D_N^{w,1})_1=(Z_w)_1
                  \quad(w\ge N).
 \tag{D-27c}
\]
It follows that \(\pi_1(D_N^{w,1})\) is a quotient of
\(\pi_1(Z_w)=k\); it is obtained by adding to the latter
the boundaries of the additional two-simplices \(a\,df\).
This also proves \(\pi_0=0\).

If \(N\le w<2N\), no product \(a\,df\) can have both
\(a\) and \(f\) of weight at least \(N\). Every coefficient
\(a\) is then in \(S\), so \(D_N^{w,1}=Z_w\). This proves
the middle range of (D-27a), and the range \(w<N\) is zero
by the definition.

Suppose now that \(w\ge2N\). Write \(a=N\), \(b=w-N\),
so \(a,b\ge N\). In simplicial degree two use cumulative
coordinates \(T,X,Y\), whose three faces send
\[
 (T,X,Y)\longmapsto(T,T,X),\quad(T,X,X),\quad(T,X,0).
 \tag{D-27d}
\]
For \(0\le s<a\), the form
\[
 \omega_s=T^{[s]}X^{[a-s]}\,d(Y^{[b]})
                 \in(D_N^{w,1})_2
\]
is allowed: its coefficient and its differentiated
function have weights \(a,b\ge N\).

Identify \(\pi_1(Z_w)\) with \(k\) by the connecting
map of (D-27b), choosing \(d(X^{[w]})\) as generator.
In particular the class represented by
\(T^{[c]}X^{[w-c-1]}dX\) is
\(\binom wc\). Using (D-27d) and the divided-power
multiplication formula, the alternating face of
\(\omega_s\) has class
\[
 \begin{aligned}
 I_s
 &=\binom as\binom wa
   -\binom ws\binom{w-s-1}{a-s}\\
 &=\binom ws\binom{w-s-1}{a-s-1}.
 \end{aligned}
 \tag{D-27e}
\]
The first equality includes the coefficient from merging
\(T^{[s]}\) with \(X^{[a-s]}\) on the first face;
omitting this coefficient would give the wrong obstruction.

These integer coefficients generate the unit ideal, because
\[
 \sum_{s=0}^{a-1}(-1)^s
       \binom ws\binom{w-s-1}{a-s-1}
             =(-1)^{a-1}.
 \tag{D-27f}
\]
For a direct proof, extract the coefficient of \(x^{a-1}\)
in the formal power-series identity
\[
 (1+x)^{w-1}
     \left(1-\frac{x}{1+x}\right)^w=(1+x)^{-1}.
\]
Only terms with \(s\le a-1\) contribute to that coefficient,
and they are exactly the summands in (D-27f).
Thus the additional boundaries kill the generator of
\(\pi_1(Z_w)\) over every \(k\). Together with (D-27c),
this proves the remaining range of (D-27a).

Finally, the vertical kernel in degree zero is \(S\):
this is immediate below weight \(N\), and is the usual
Poincar\'e kernel above it. In weights at least \(N\),
every vertically closed element of \(D_N^1\) is closed
in \(F^1\), hence is \(df\) for \(f\in F^0=D_N^0\)
in that weight. Below weight \(N\) there are no one-forms.
This proves the assertion about vertical degrees zero and one.
\end{proof}

\paragraph{The image remains genuinely nonfree.}
The preceding result does not arise by eventual equality
\(D_N=F\). For example, take \(k=\Z\), \(N=2\), set the
relative coordinate \(T\) to zero, and consider the
two-variable multidegree \((n,1)\), \(n\ge3\). Put
\[
 U=X^{[n-1]}Y\,dX,\qquad
 V=X^{[n]}\,dY,\qquad
 W=X^{[n-1]}\,dX\,dY.
 \]
If
\[
 d_n=\gcd\left\{\binom{n-1}{r}:1\le r\le n-2\right\},
\]
the form-one and form-two image lattices in this
multidegree are respectively
\[
 \Z(U+V)+\Z d_nV,\qquad \Z d_nW.
 \tag{D-27g}
\]
To see the first formula, differentiate \(X^{[n]}Y\)
and then list the monomial products \(a\,df\), with
both nonconstant factors of weight at least two.
Modulo \(U+V\), their coefficients are exactly the
internal binomial coefficients in the definition of \(d_n\).
For the second, products \(df\,dg\) already give their
gcd. Any additional coefficient factor gives a
multinomial coefficient divisible by the same gcd:
the only nontrivial case is \(n-1=\ell^r\), where
reduction of \((x+y+z)^{\ell^r}\) modulo \(\ell\)
shows the divisibility. The binomial gcd is \(\ell\)
in that case and \(1\) otherwise, as proved in
Statement~\ref{sa08:C11}.
Hence the image has a nontrivial lattice index in
unbounded weights. Nevertheless its vertical
multidegree complex is exact, since
\[
 d(X^{[n]}Y)=U+V,\qquad dV=W,\qquad dU=-W.
\]

\paragraph{Finite checks and the remaining gate.}
We record the executed programs and their precise scope.

The program \texttt{code/D\_tail\_differential\_image.py}
checks the actual alternating face images in 264 cases:
\(N=2,3,4\), \(p=2,3,5,7\), and \(N\le w\le24\).
They agree with (D-27a).
The separate program
\texttt{code/D\_tail\_differential\_H2.py} finds
\(\pi_2(D_N^{w,1})=0\) in 45 cases over
\(\F_2,\F_3,\F_5\), for \(N=2,3\) and \(N\le w\le9\).
These field computations concern the image defined over
that field; an image need not commute with reduction
of the integral lattice.
Exact integral multidegree calculations in
\texttt{code/D\_tail\_vertical\_multidegree.py} and
\texttt{code/D\_tail\_vertical\_threevars.py} find no
positive vertical cohomology in respectively 194
two-variable and 108 three-variable tested cases.
Their matrices, counts, and outputs are saved in
the correspondingly named log files.

Thus actual nonfree differential relations remove the
specific first-horizontal-homotopy obstruction detected
for a free coefficient resolution in
Statement~\ref{sa08:C15}. The calculations do not prove
eventual acyclicity in every horizontal dimension and
every fixed form degree, nor do they construct a
comparison source for arbitrary targets. In particular
the degree-zero standard tail still need not map to an
arbitrary theory, by Statement~\ref{sa08:C13}.
